\documentclass[a4paper]{article}

\usepackage[english, russian]{babel}
\usepackage[utf8]{inputenc}
\usepackage{amsmath,amssymb,amsthm,mathtools}
\usepackage{graphicx}
\usepackage[ruled,linesnumbered,vlined]{algorithm2e}
\usepackage{comment}
\usepackage{hyperref}
\usepackage{epigraph}

% Margins
\usepackage[top=2.5cm, left=3cm, right=3cm, bottom=3.0cm]{geometry}
% Colour table cells
\usepackage[table]{xcolor}
\definecolor{cadmiumgreen}{rgb}{0.0, 0.42, 0.24}
\definecolor{pigmentblue}{rgb}{0.2, 0.2, 0.6}
\hypersetup{ %use colored text instead of ugly boxes
	colorlinks,
	linkcolor={cadmiumgreen},
	urlcolor={pigmentblue}
}

\usepackage{tikz}
\usepackage{pdfpages}

% My commands
\newcommand{\statement}[2]{\textit{\textbf{#1} #2}}
\newcommand{\e}{\varepsilon}
\newcommand{\needpicture} {
    {
        \begin{center}
        \rowcolors{1}{yellow!75}{yellow!75}
        \begin{tabular}{|c|}
            \hline
            Здесь нужна картинка\\
            \hline
        \end{tabular}
        \end{center}
    }
}
\newcommand{\heart}{\ensuremath\heartsuit}

\newtheorem{theorem}{Теорема}[section]
\newtheorem{definition}{Определение}[section]
\newtheorem{example}{Пример}[section]
\newtheorem{problem}{Задача}[section]
\newtheorem{lemma}{Лемма}[section]
\newtheorem{corollary}{Следствие}[section]
\newtheorem{nb}{NB}[section]

\newtheorem*{theorem*}{Теорема}
\newtheorem*{definition*}{Определение}
\newtheorem*{example*}{Пример}
\newtheorem*{problem*}{Задача}
\newtheorem*{lemma*}{Лемма}
\newtheorem*{nb*}{NB}

% My environments
\newcommand*{\QEDTRIANGLE}{\null\nobreak\hfill\ensuremath{\bigtriangleup}} 
\newenvironment{algodescription}[1]
    {\noindent\textbf{#1.}}
    {\QEDTRIANGLE}

%Author tags
\newcommand{\groth} {Нечаев~Е.}

%%%%%%%%%%%%%%%%%
%     Title     %
%%%%%%%%%%%%%%%%%
\title{Алгосы, часть ii}
\author{Денис Осипов, Иван Ермошин, Егор Нечаев}
\date{\today}

\begin{document}

%\includepdf{frontcover.pdf}

\maketitle

\section*{Введение}

Этот проект~-- коллективный \textbf{конспект} по второй части курса <<Математические основы алгоритмов>>, впервые прочитанного первокурсникам МКН СПбГУ в первой половине ii семестра 2020 года Эдуардом Алексеевичем Гиршем.

Актуальные исходники: \url{https://www.overleaf.com/read/hnbkrkyknbpk} и \url{https://github.com/gogochushij/algosi-hirsch}

Если вы хотите \textbf{принять участие} в написании билетов, или же
\textbf{сообщить об ошибке}, напишите \url{http://vk.com/gogochushij}. Предполагается, что каждый автор напишет около 4 билетов, но мы будем рады любой посильной помощи. \href{https://docs.google.com/spreadsheets/d/17MKhLVzCyYvEKlm6W5Bb-6uUDNyv0QLmdvuh4N6JfXI/edit?usp=sharing}{\texttt{Здесь}} можно посмотреть, с какими билетами вы можете помочь проекту.

Символом \heart~обозначены разделы, которые не входят в экзамен (еще не всё проставлено).

\tableofcontents \newpage

%%%%%%%%%%%%%%%%%%%%%%%%%%%%%%%%%%%%%%%%%
% Здесь настроить конспект для компиляции
% Чтобы скомпилировать все, что есть, впишите:
% \foreach \n \in {1,...,15,16,17,18,19,20,21} {\input{que\n.tex}} 
% ВМЕСТО строки ниже
\foreach \n in {1,...,15,16,17,18,19,20,21} {\input{que\n.tex}}
%%%%%%%%%%%%%%%%%%%%%%%%%%%%%%%%%%%%%%%%%

% que6.tex ссылается на que5.tex
% que21.tex ссылается на que20.tex
% que9.tex -> que8.tex -> que7.tex

\end{document}
